%========================================================
% SECTION 7.3 — Global Depolarizing Channel
%========================================================

\newpage

\section{Experiment 3: Global Depolarizing Channel}
\label{sec:global_depolarizing_channel}

%--------------------------------------------------------
% Subsection 7.3.1 — Objective and Physical Motivation
%--------------------------------------------------------

\subsection*{Objective and Physical Motivation}

This experiment investigates the degradation of a maximally entangled Bell state under an effective two-qubit depolarizing channel.

Unlike the previous experiments, which were based on coherent unitary evolution or statistical averaging over unitary realizations, the present model introduces a genuinely non-unitary transformation acting directly at the level of the density operator.

The objective is to characterize how depolarization affects:

\begin{itemize}
    \item purity and mixedness;
    \item fidelity with respect to the Bell reference state;
    \item concurrence and entanglement degradation;
    \item observable correlations;
    \item the structure of the correlation tensor.
\end{itemize}

This experiment therefore establishes the first explicit decoherence model of the simulator and provides a baseline description of irreversible coherence loss in bipartite quantum propagation.

%--------------------------------------------------------
% Subsection 7.3.2 — Model and Analytical Predictions
%--------------------------------------------------------

\subsection*{Model and Analytical Predictions}

The initial state is the Bell state

\[
|\Psi^+\rangle
=
\frac{|01\rangle + |10\rangle}{\sqrt{2}},
\qquad
\rho_0
=
|\Psi^+\rangle
\langle\Psi^+|.
\]

The depolarizing channel is defined as

\begin{equation}
\rho(p)
=
(1-p)\rho_0
+
\frac{p}{4}I_4,
\qquad
p \in [0,1].
\label{eq:depolarizing_channel_experiment}
\end{equation}

The parameter \(p\) controls the strength of the depolarization process.

The limiting cases are

\[
\rho(0)
=
\rho_0,
\qquad
\rho(1)
=
\frac{I_4}{4}.
\]

Thus, the channel continuously interpolates between a pure maximally entangled Bell state and the maximally mixed two-qubit state.

\medskip

For this family of states, the analytical predictions are

\[
\gamma_{AB}(p)
=
1
-
\frac{3}{2}p
+
\frac{3}{4}p^2,
\qquad
\gamma_A
=
\gamma_B
=
\frac{1}{2},
\]

\[
F_{\Psi^+}(p)
=
1
-
\frac{3}{4}p,
\qquad
C(p)
=
\max
\left(
0,
1-\frac{3}{2}p
\right).
\]

The concurrence vanishes at

\[
p
=
\frac{2}{3},
\]

marking the transition from entangled to separable states.

\medskip

The action of the depolarizing channel on the correlation tensor corresponds to a uniform contraction of the Bell-state tensor:

\[
T(p)
=
(1-p)
T_{\Psi^+}.
\]

Using the Bell-state tensor

\[
T_{\Psi^+}
=
\begin{pmatrix}
1 & 0 & 0 \\
0 & 1 & 0 \\
0 & 0 & -1
\end{pmatrix},
\]

one obtains

\begin{equation}
T(p)
=
\begin{pmatrix}
1-p & 0 & 0 \\
0 & 1-p & 0 \\
0 & 0 & -(1-p)
\end{pmatrix}.
\label{eq:depolarizing_tensor}
\end{equation}

The corresponding diagonal correlations are therefore

\[
c_{xx}(p)
=
1-p,
\qquad
c_{yy}(p)
=
1-p,
\qquad
c_{zz}(p)
=
-(1-p).
\]

These expressions are consistent with the correlation-tensor formalism introduced in Section~\ref{sec:correlation_tensor_representation} and with the analytical derivations reported in Appendix~\ref{app:depolarizing_channel_analytics}.

%--------------------------------------------------------
% Subsection 7.3.3 — Numerical Results
%--------------------------------------------------------

\subsection*{Numerical Results}

The depolarization parameter \(p\) is sampled over the interval

\[
p \in [0,1],
\]

and all state-characterization quantities and correlation observables are evaluated numerically and compared with the corresponding analytical predictions.

The numerical results show agreement at machine precision for all quantities considered in the simulation.

In particular:

\begin{itemize}
    \item the Frobenius-norm error of the correlation tensor remains of order \(10^{-16}\);
    \item the errors on purity, fidelity, concurrence, and diagonal correlations remain at numerical roundoff level.
\end{itemize}

\begin{figure}[htbp]
\centering
\includegraphics[width=\linewidth]{figures/experiment_3/experiment_3_depolarizing_channel_summary.png}
\caption{
Numerical and analytical behavior of purity, fidelity, concurrence, and correlations under the depolarizing channel.
Agreement is exact up to numerical precision.
}
\label{fig:experiment3_depolarizing_summary}
\end{figure}

Figure~\ref{fig:experiment3_depolarizing_summary} summarizes the numerical evolution of the principal observables under the depolarizing channel.


The numerical and analytical curves are visually indistinguishable within numerical precision.

The global purity decreases monotonically from unity toward the maximally mixed limit, while the concurrence exhibits threshold behavior and vanishes at

\[
p
=
\frac{2}{3}.
\]

At the same time, the diagonal correlations decay linearly in magnitude, consistently with the tensor contraction law given in Eq.~\eqref{eq:depolarizing_tensor}.

%--------------------------------------------------------
% Subsection 7.3.4 — Discussion
%--------------------------------------------------------

\subsection*{Discussion}

This experiment validates the first explicit non-unitary extension of the simulator.

Unlike the fixed-phase and ensemble-based phase models studied previously, the depolarizing channel introduces an isotropic loss of coherence directly at the level of the two-qubit density operator.

The numerical results show that the different observables exhibit distinct degradation behaviors:

\begin{itemize}
    \item the global purity decreases smoothly over the full interval;
    \item the concurrence exhibits a sharp entanglement threshold;
    \item the observable correlations decay linearly;
    \item the full correlation tensor undergoes a uniform isotropic contraction.
\end{itemize}

The reduced density operators remain maximally mixed throughout the evolution:

\[
\rho_A
=
\rho_B
=
\frac{I}{2}.
\]

This behavior is consistent with the results of Experiment~\ref{sec:exp_2_random_phase_ensemble_sweep}, 
where local mixedness also remained constant despite degradation of global coherence.

However, the physical origin of mixedness is fundamentally different in the two cases.

In the random-phase ensemble model, mixedness arises from averaging over unresolved coherent realizations while preserving an underlying phase structure.

In contrast, the depolarizing channel suppresses the nontrivial two-qubit correlations isotropically by mixing the state directly with the maximally mixed operator \(I_4/4\).

The experiment therefore establishes a clear distinction between:

\begin{itemize}
    \item decoherence emerging from statistical phase uncertainty;
    \item irreversible isotropic depolarization modeled through a CPTP channel.
\end{itemize}

%--------------------------------------------------------
% Subsection 7.3.5 — Physical Interpretation
%--------------------------------------------------------

\subsection*{Physical Interpretation}

The depolarizing channel considered in this experiment is a global two-qubit depolarizing channel. It acts directly on the complete bipartite density operator and models the progressive replacement of the initial Bell state by the maximally mixed two-qubit state,

In this sense, the parameter \( p \) does not represent the degradation of a specific fiber arm. Instead, it controls an isotropic loss of information at the level of the full two-photon state. The limiting cases are directly interpretable: for \( p=0 \), the initial Bell state is recovered, while for \( p=1 \), the output state is completely mixed.

\medskip

Physically, this model should be interpreted as an effective baseline for incoherent two-qubit degradation rather than as a microscopic description of optical-fiber propagation. It captures the idealized situation in which all nontrivial information contained in the bipartite state is uniformly suppressed, without distinguishing between different physical origins or between the two propagation arms.

\medskip

Under global depolarization, all traceless components of the two-qubit state are contracted uniformly. As a consequence, the observable two-photon correlations, the fidelity with respect to the reference Bell state, and the concurrence decrease monotonically as \( p \) increases. In the fully depolarized limit, the state \( I_4/4 \) contains no preferred polarization structure, no entanglement, and no nonzero two-qubit correlation tensor components.

\medskip

This behavior differs from the random-phase ensemble studied previously. In the random-phase model, decoherence arises from statistical averaging over unresolved phase realizations and therefore preserves a specific phase-related structure. By contrast, the global depolarizing channel removes information isotropically: no particular phase direction, measurement basis, or fiber arm is privileged.

\medskip

For this reason, the global depolarizing channel is not intended to be the final physical model of fiber propagation. Its role is to provide a simple and analytically transparent reference case for mixed-state degradation. More realistic descriptions, such as local depolarization acting independently on arms \( A \) and \( B \), will refine this baseline by assigning the loss mechanism to the individual propagation paths.